% TODO(team): Keep targeted availability separate from buffered playback and permanent loss.
\section{E4: second-holder failure}
\begin{frame}{Can the replacement serve after both holders stop?}
\begin{enumerate}
\item Stop the first original holder.
\item Wait a fixed \textbf{4 / 14 / 30 seconds}, independent of repair success.
\item Stop the second original holder.
\item Probe a preselected late-queue segment from its replacement.
\end{enumerate}
\small S / F50 / L; two paired contexts per delay. Success requires correct hash \emph{and} committed registration.
\takeaway{A buffered viewer may finish even when fresh delivery is unavailable.}
\note{The probe checks one sentinel, not the whole video. Its late queue position makes repair timing important.}
\end{frame}
\begin{frame}{The selected replacement was not ready in time}
\begin{tabularx}{\textwidth}{lYYY}
\toprule Policy & Second fault +4 s & +14 s & +30 s \\\midrule
S & 0/2 & 0/2 & 0/2\\
F50 & 0/2 & 0/2 & 0/2\\
L & 0/2 & 0/2 & 0/2\\\bottomrule
\end{tabularx}
\vspace{4mm}
All 18 probes were unavailable under the specified success condition.
\takeaway{A boundary was observed; no policy advantage or permanent data loss was established.}
\note{Stopped processes retained their disks. Later delays were not tested. A future completed-repair control should verify replacement availability after every required segment is repaired. Source: compact-study-results.md, E4.}
\end{frame}
